% TauCetiNameMap.tex -- proposed TauCeti names for the foundational layer of lean-DLMF.
% Compile with LuaLaTeX: latexmk -lualatex TauCetiNameMap.tex
\documentclass[11pt,a4paper]{article}
\usepackage{leanreport}
\hypersetup{pdftitle={lean-DLMF: TauCeti name map of the foundational layer}}

\begin{document}

\title{TauCeti Name Map\\of the Foundational Layer}
\author{Bastiaan J. Braams}
\date{2026-11-04}
\maketitle
\sloppy

\begin{abstract}
  The foundational (\texttt{Basic}) layer of \texttt{lean-DLMF} is to be offered to TauCeti
  through roadmaps. In TauCeti, \texttt{lean-DLMF} declarations get the namespace
  \lean{TauCeti}, and a later move to Mathlib drops that prefix. Every foundational name
  therefore has to be a good Mathlib name once the \lean{DLMF.\textless Area\textgreater}
  prefix is removed. This report applies a fixed mapping rule to all 3,660 public
  foundational declarations, finds the clashes with Mathlib, TauCeti and lean-codes and the
  clashes between areas, and proposes renames for the abbreviated names. The complete table
  is \file{reports/TauCetiNameMap.tsv}. The report closes with the decisions needed before
  the renaming round, as adopted.
\end{abstract}

\section{Mapping rule}

The project keeps its area namespaces. A roadmap maps names as follows.
\begin{itemize}
  \item \lean{DLMF.\textless Area\textgreater.x} becomes \lean{TauCeti.x} (root placement) or
    \lean{TauCeti.Complex.x} (complex placement). Root placement is the default. Complex
    placement is used only where root placement would clash with an upstream name, which in
    practice means that Mathlib or TauCeti has a real-valued function of that name at the root.
  \item \lean{DLMF.\textless Area\textgreater.Real.x} becomes \lean{TauCeti.Real.x}.
  \item A lemma in the namespace of a definition or structure keeps that namespace
    (\lean{DLMF.Airy.IsAirySol.eq\_zero} becomes \lean{TauCeti.IsAirySolution.eq\_zero}).
  \item A dot-notation lemma outside \lean{DLMF} keeps its name under \lean{TauCeti}.
  \item The general areas (Chapter~1 \lean{AlgebraicAnalytic}, Chapter~2
    \lean{AsymptoticApproximations} and \lean{DLMF.Core}) are not special functions. Their
    roadmaps place declarations in the namespace of the Mathlib API they extend, which is
    chosen per file. The table still gives the root form, marked \texttt{general}.
\end{itemize}

\section{The table}

\file{reports/TauCetiNameMap.tsv} has one row per public declaration of a foundational
module. Generated auxiliaries are excluded: recursors, \lean{match\_} and \lean{proof\_}
terms, and \lean{eq\_1}. The columns are:
\begin{enumerate}
  \item kind: \texttt{def}, \texttt{type} or \texttt{thm};
  \item the current name;
  \item the module;
  \item the proposed new \texttt{lean-DLMF} name, or \texttt{=} if unchanged;
  \item the proposed TauCeti name;
  \item the placement: \texttt{root}, \texttt{Complex}, \texttt{Real}, \texttt{dot API}
    (namespace of a definition or structure), \texttt{sub-namespace} (helper namespace),
    \texttt{general} or \texttt{dot};
  \item flags: upstream clash, target clash, possible duplicate, remaining abbreviation.
\end{enumerate}
The table has 1,624 definitions, 53 types and 1,983 theorems. The placement counts are
2,783 root, 361 general, 311 dot API, 151 Real, 28 sub-namespace, 16 Complex, 9 upstream
(replaced by the upstream function or stated about it) and 1 dot. 732 names have a rename
proposal, 370 of them definitions, types or namespaces of definitions.
Theorem names follow the definitions they mention: \lean{pcfU\_neg} becomes
\lean{parabolicCylinderU\_neg}.

Renaming the foundational definitions also renames about 2,200 theorem-layer declarations
whose names contain them, and changes about 10,000 identifier occurrences across the
library. The renaming round covers both layers. Every rename is recorded in
\file{reports/RenameMap20261101.tsv}.

\section{Clashes}

\subsection{Upstream clashes}

Under root placement, eight foundational declarations clash with an upstream declaration.
\begin{itemize}
  \item \lean{Confluent.lowerIncGamma}, proposed \lean{lowerIncompleteGamma}: TauCeti has the
    real \lean{TauCeti.lowerIncompleteGamma}. The same holds for
    \lean{Statistical.normIncBetaI}, proposed \lean{regularizedIncompleteBeta}, against
    TauCeti's real \lean{regularizedIncompleteBeta}, and for
    \lean{MatrixArgument.multivariateGamma} with its lemma \lean{multivariateGamma\_one}. The
    bridges already exist. Each whole family goes in \lean{Complex}, including the
    \lean{OnLog} variants, \lean{upperIncompleteGamma}, the complex \lean{incompleteBeta} and
    \lean{multivariateGammaVec}. The complex names mirror TauCeti's real ones, so that the
    bridges read \lean{Complex.X\_ofReal}: in particular $P(a,z)$ becomes
    \lean{Complex.regularizedGamma} after TauCeti's \lean{regularizedGamma}. The function
    \lean{Statistical.Real.normIncBetaI} (now \lean{Real.regularizedIncompleteBeta}) has complex
    parameters, a real argument and complex values; it agrees with TauCeti's real
    \lean{regularizedIncompleteBeta} only for real positive parameters, so it gets a bridge
    there and is not replaced. \lean{Statistical.normIncGammaP} and \lean{normIncGammaQ}
    (quotient forms) agree with the \lean{Confluent} continuations except at the poles of
    $\Gamma(a)$, where the quotients are totalized to~$0$; keeping one of them needs a proof
    round, because every use must be checked at those points.
  \item \lean{AlgebraicAnalytic.hermiteFunction} against TauCeti's real
    \lean{hermiteFunction}. By 18.7.11 DLMF's
    $(\sqrt\pi\,2^nn!)^{-1/2}e^{-x^2/2}H_n(x)$ equals TauCeti's
    $\mathit{He}_n(\sqrt2x)e^{-x^2/2}/\sqrt{n!\sqrt\pi}$; the local definition differs only by
    the cast to $\mathbb C$. It is replaced by the TauCeti function, and the tag 1.18.E42 moves
    to a theorem about its cast.
  \item \lean{AlgebraicAnalytic.dlmfResolventCubic}: TauCeti's \lean{resolventCubic} is the
    Ferrari cubic $X^3-pX^2-4rX+(4pr-q^2)$, whereas DLMF~1.11.18 is
    $X^3-2pX^2+(p^2-4r)X+q^2$, with roots $-(\alpha_1+\alpha_{j+1})^2$. The two are different
    polynomials. The DLMF one becomes \lean{eulerResolventCubic}, after Euler's method.
  \item \lean{OrthogonalPolynomials.poissonKernel} against Mathlib's \lean{poissonKernel}
    (the disc Poisson kernel). The DLMF object is the kernel
    $\sum p_n(x)p_n(y)z^n/h_n$ of 18.2.41. Proposal: \lean{orthogonalPoissonKernel}, at the
    root.
  \item \lean{Hypergeometric.ascPochhammer\_eval\_add} clashes with the lean-codes lemma of
    the same name in the root namespace (module \module{Pochhammer}). If lean-codes is
    offered upstream too, one of the two must go.
\end{itemize}
There are no clashes with Mathlib or TauCeti in the \lean{Real} namespace. The lean-codes
module \module{ToMathlib.Analysis.SpecialFunctions.Bessel} defines \lean{Complex.besselI}
and related lemmas; \lean{DLMF.Bessel.besselI} at the root does not clash with them, but a
TauCeti submission of both would need one canonical modified Bessel function.

\subsection{Clashes between areas}

Two pairs of foundational definitions mapped to the same TauCeti name. Both were duplicates
and are merged:
\begin{itemize}
  \item \lean{AlgebraicAnalytic.stieltjesTransform} (1.14.47) is the only definition of the
    Stieltjes transform; DLMF 2.6.8 is the theorem
    \lean{AsymptoticApproximations.stieltjesTransform\_eq}.
  \item \lean{AlgebraicAnalytic.IsLinearODE2SolutionOn} (1.13.1) is the only second-order
    linear ODE predicate. It allows an inhomogeneous term; DLMF 2.7.1 is the homogeneous case,
    stated as \lean{AsymptoticApproximations.isLinearODE2SolutionOn\_iff\_isODE2SolutionOn},
    and \lean{Heun.IsLinearODE2On} is an abbreviation of the homogeneous case with the API
    \lean{differentiableAt}, \lean{hasDerivAt\_deriv} and \lean{mk}.
\end{itemize}
Three further pairs would clash under root placement:
\begin{itemize}
  \item \lean{AlgebraicAnalytic.mobius} (1.9.40, a function of four complex
    coefficients) and \lean{Weierstrass.mobius} (23.15.3, an action of
    $\mathrm{SL}(2,\mathbb Z)$): \lean{mobiusMap} and \lean{modularMobius}.
  \item \lean{AlgebraicAnalytic.laplaceTransform} and \lean{MatrixArgument.laplaceTransform}:
    the matrix one becomes \lean{matrixLaplaceTransform}.
  \item \lean{AlgebraicAnalytic.laurentCoeff} ($a_n$ about $z_0$) is the only Laurent
    coefficient; the Chapter~2 coefficient about $0$ is the theorem
    \lean{AsymptoticApproximations.laurentCoeff\_zero\_eq} (2.10.26).
\end{itemize}
The segment integral over $[0,w]$ of Chapter~19 is \lean{Core.primitive}, with the unfolding
lemma \lean{EllipticIntegrals.primitive\_eq\_integral}.

\subsection{Further possible duplicates}

\begin{itemize}
  \item \lean{EllipticIntegrals.compEllIntKpk} and \lean{EllipticIntegrals.cCompEllIntKk} are
    both $K'(k)$, continued from different domains ($k^2\notin(-\infty,0]$ and
    $k\notin(-\infty,0]$).
  \item The forward difference of Chapter~2 is Mathlib's \lean{fwdDiff 1}.
  \item \lean{AlgebraicAnalytic.fourierTransform} uses DLMF's normalization
    $(2\pi)^{-1/2}\int f(t)e^{ixt}\,dt$. Mathlib's \lean{Real.fourierIntegral} uses
    $e^{-2\pi i xt}$, so the two are different functions; the bridge belongs in the
    Chapter~1 roadmap.
  \item \lean{BernoulliEuler.binomialSeriesCoeff} (a polynomial over $\mathbb Q$ in the
    exponent) and \lean{AsymptoticApproximations.binomSeriesCoeff} (its value at a complex
    exponent, for complex coefficients) are both kept; for rational coefficient ratios
    \lean{AsymptoticApproximations.binomSeriesCoeff\_eq\_eval\_binomialSeriesCoeff} identifies
    them.
\end{itemize}

\section{Rename proposals}

The proposals spell out abbreviations, follow upstream precedent (TauCeti's
\lean{lowerIncompleteGamma}, \lean{regularizedIncompleteBeta}, \lean{logIntegral}; Mathlib's
\lean{fwdDiff}) and keep DLMF's own function letters (\lean{whittakerM}, \lean{ferrersP},
\lean{mathieuCe}, \lean{spheroidalPs}). The families are:
\begin{itemize}
  \item Elliptic integrals, Chapter~19: \lean{compEllIntKk} $\to$ \lean{completeEllipticK},
    likewise \lean{E}, \lean{Pi}, \lean{D}; \lean{compEllIntKpk} $\to$
    \lean{completeEllipticKPrime}; \lean{cCompEllIntKk} $\to$ \lean{complementaryEllipticK};
    \lean{incEllIntFk} $\to$ \lean{incompleteEllipticF}, likewise \lean{E}, \lean{Pi},
    \lean{D}; \lean{symEllIntRF} $\to$ \lean{symmetricEllipticRF}, likewise \lean{RC},
    \lean{RD}, \lean{RG}, \lean{RJ} and the general \lean{symEllIntR} $\to$
    \lean{symmetricEllipticR}; \lean{symEllIntS} $\to$ \lean{carlsonRootProduct};
    \lean{SymEllIntArg} $\to$ \lean{IsCarlsonArgument}; \lean{cauchyPVIoi} $\to$
    \lean{principalValueIoi}.
  \item Parameter forms of the older \lean{Elliptic} area (argument $m=k^2$):
    \lean{compEllIntKm} $\to$ \lean{completeEllipticKOfParameter}, likewise for $E$, $\Pi$,
    $F$, $\mathrm{am}$, $\mathrm{sn}$, $\mathrm{cn}$, $\mathrm{dn}$ and arcsn (decision~1).
  \item Jacobi elliptic functions, Chapter~22: the suffix \lean{k} is dropped
    (\lean{jacobiSnk} $\to$ \lean{jacobiSn}, for all twelve letter pairs, and
    \lean{jacobiNomek}, \lean{jacobiComodulusk}, \lean{jacobiEpsilonk}, \lean{jacobiTauk},
    \lean{jacobiPQk}). The $\tau$ forms become \lean{\ldots OfTau} (\lean{jacobiSnTau}
    $\to$ \lean{jacobiSnOfTau}, \lean{compEllIntKTau} $\to$ \lean{completeEllipticKOfTau},
    \lean{jacobiNevilleTau} $\to$ \lean{nevilleThetaOfTau}). The suffix \lean{m} of
    \lean{JacobiElliptic.Real.jacobiArcndm} and its siblings is dropped, because they take
    the modulus $k$. Two names say something else than they mean: \lean{jacobiZetak} is the
    theta argument $\zeta=\pi z/(2K)$ of 22.2.3, proposed \lean{jacobiThetaArgument}, and
    \lean{jacobiZetaFunctionk} is Jacobi's $\mathrm{Z}$ function, proposed \lean{jacobiZeta}.
  \item Parabolic cylinder functions, Chapter~12: \lean{pcfX} $\to$
    \lean{parabolicCylinderX} for all eighteen definitions (\lean{U}, \lean{V}, \lean{D},
    \lean{UBar}, \lean{Even}, \lean{Odd}, the moduli \lean{F}, \lean{G}, the phases
    \lean{Theta}, \lean{Psi}, \lean{Moment} and the derivatives and values at~$0$).
  \item Chapters~6--8: \lean{sinInt} $\to$ \lean{sinIntegral}, likewise \lean{cosInt},
    \lean{sinhInt}, \lean{coshInt}, \lean{cosIntCin}, \lean{sinIntTail}; \lean{expIntEin}
    $\to$ \lean{expIntegralEin}; \lean{expIntEi} $\to$ \lean{expIntegralEi};
    \lean{logIntLi} $\to$ \lean{logIntegralLi} (TauCeti's \lean{logIntegral} starts at~$2$);
    \lean{genExpIntE} $\to$ \lean{expIntegralE}; \lean{fresnelCosInt} $\to$ \lean{fresnelC},
    \lean{fresnelSinInt} $\to$ \lean{fresnelS}; \lean{expIntAuxF} $\to$
    \lean{sinIntegralAuxiliaryF}, \lean{fresnelAuxF} $\to$ \lean{fresnelAuxiliaryF};
    \lean{lowerIncGamma} $\to$ \lean{lowerIncompleteGamma}, \lean{upperIncGamma} $\to$
    \lean{upperIncompleteGamma}, \lean{normIncGammaPContinued} $\to$
    \lean{regularizedGamma}; \lean{incBeta} $\to$ \lean{incompleteBeta},
    \lean{normIncBetaI} $\to$ \lean{regularizedIncompleteBeta}; \lean{scIncGamma} $\to$
    \lean{tricomiIncompleteGamma}.
  \item Generalized hypergeometric functions, Chapter~16: \lean{genHG} $\to$
    \lean{generalizedHypergeometric}, \lean{genRegularizedHG} $\to$
    \lean{regularizedGeneralizedHypergeometric}, \lean{genCoeff} $\to$
    \lean{generalizedHypergeometricCoeff}, \lean{genPoch} $\to$ \lean{pochhammerProduct},
    \lean{genGamma} $\to$ \lean{gammaProduct}, \lean{genIntPoch} $\to$
    \lean{pochhammerOfInt}, \lean{genBilateralHG} $\to$ \lean{bilateralHypergeometric}; the
    predicates \lean{GenWellPoised} $\to$ \lean{IsWellPoised}, \lean{GenVeryWellPoised},
    \lean{GenKBalanced}, \lean{GenTerminates} $\to$ \lean{IsTerminating},
    \lean{GenDenRegular} $\to$ \lean{LowerParametersRegular}; \lean{poch} $\to$
    \lean{pochhammer}; \lean{regularizedGaussHG} $\to$ \lean{regularizedGaussHypergeometric},
    \lean{gaussHGCoeff} $\to$ \lean{gaussHypergeometricCoeff}.
  \item Appell functions: \lean{appF1} $\to$ \lean{appellF1}, and so on for all four
    functions with their \lean{Coeff}, \lean{Euler} and \lean{EulerKernel} companions.
  \item Matrix argument, Chapter~35: \lean{hyp0F1} $\to$ \lean{matrixHypergeometric0F1} and
    the other five, \lean{hypergeometricPFQ} $\to$ \lean{matrixHypergeometricPFQ},
    \lean{jacobiP} $\to$ \lean{matrixJacobiP}, \lean{laguerreL} $\to$ \lean{matrixLaguerreL},
    \lean{besselA}, \lean{besselB}, \lean{kummerPsi}, \lean{laplaceTransform} likewise with the
    prefix \lean{matrix}; \lean{etr} $\to$ \lean{expTrace}. Without the prefix these names
    would be read as scalar functions at the TauCeti root.
  \item Lamé, Chapter~29: the modulus suffix is dropped (\lean{lameEck} $\to$ \lean{lameEc},
    \lean{lameAk} $\to$ \lean{lameA}, \lean{lameUEk} $\to$ \lean{lameUE}, and the other
    polynomial forms); \lean{Poly}, \lean{Rec}, \lean{CF} are spelled out.
  \item Heun, Chapter~31: \lean{heunHypK} $\to$ \lean{heunHypergeometricK} and the rest of
    the 31.11 family; \lean{heunRecP} $\to$ \lean{heunRecurrenceP}; \lean{heunCFTail} $\to$
    \lean{heunContinuedFractionTail}; \lean{heunEps} $\to$ \lean{heunEpsilon};
    \lean{IsHeunSolOn} $\to$ \lean{IsHeunSolutionOn} and the confluent forms.
  \item Spheroidal, Chapter~30: the continued-fraction coefficients of 30.3.6 and 30.3.7
    become \lean{spheroidalPowerAlpha} and \lean{spheroidalLegendreAlpha} (and the
    \lean{Beta}, \lean{Gamma} versions) in place of \lean{spheroidalCFAlpha} and
    \lean{spheroidalCFAlphaAlt}; \lean{spheroidalEll} $\to$
    \lean{spheroidalEigenvalueCoeff}; \lean{spheroidalRecA} $\to$
    \lean{spheroidalRecurrenceA}.
  \item Coulomb, Chapter~33: the functions $f,h,s,c$ of $\S$33.14 are \lean{coulombf},
    \lean{coulombh}, \lean{coulombs}, \lean{coulombc}, which differ from \lean{coulombF} only
    in case. They become \lean{coulombEnergyF}, \lean{coulombEnergyH},
    \lean{coulombEnergyS}, \lean{coulombEnergyC} (decision~4).
  \item Names that contain ``DLMF'': \lean{gaussSumDLMF} $\to$ \lean{dirichletGaussSum},
    \lean{IsPrimitiveDLMF} $\to$ \lean{IsPrimitiveByDivisors} (with a bridge to Mathlib's
    \lean{DirichletCharacter.IsPrimitive}), \lean{discriminantDLMF} $\to$
    \lean{modularDiscriminant}, \lean{ProdConvergesDLMF} $\to$
    \lean{InfiniteProductConverges}, \lean{dlmfD} $\to$ \lean{normalizedPartialDeriv},
    \lean{dlmfFourierSchwartz} $\to$ \lean{fourierTransformSchwartz},
    \lean{dlmfResolventCubic} $\to$ \lean{eulerResolventCubic}.
  \item Other abbreviations: \lean{Sol} $\to$ \lean{Solution} (\lean{IsAirySol},
    \lean{IsWeberSol}, \lean{IsSchrodingerSolOn}, \ldots); \lean{Op} $\to$ \lean{Operator}
    (\lean{besselOp}, \lean{opSpectrum}, \ldots); \lean{Asymp} $\to$ \lean{Asymptotic};
    \lean{cfTail} $\to$ \lean{continuedFractionTail}; \lean{powC} $\to$
    \lean{cpowArithmeticFunction}, \lean{sigmaC} $\to$ \lean{sigmaCpow}; \lean{factQ} $\to$
    \lean{ratFactorial}; \lean{laplacianMv} $\to$ \lean{laplacianEuclidean};
    \lean{envAi} $\to$ \lean{envelopeAi}; \lean{lgApproximant} $\to$
    \lean{liouvilleGreenApproximant}; \lean{legendrePolyP} $\to$ \lean{legendrePolynomialP}
    and the other \lean{Poly} names; \lean{basicHG}, \lean{terminatingHG} $\to$
    \lean{basicHypergeometric}, \lean{terminatingHypergeometric}.
  \item Real forms with a \lean{Real} suffix: \lean{Airy.airyAiReal} and three siblings move
    to \lean{DLMF.Airy.Real.airyAi}; likewise \lean{invPowScaleReal} and
    \lean{IsSpheroidalSolutionReal}.
  \item Selberg helpers (\lean{vdm}, \lean{iseg}, \lean{cube}, \lean{mono}, \lean{reflect},
    \lean{succPre}, \lean{ibpF}, \ldots): spelled out (\lean{vandermondeAbsPow},
    \lean{initialSegment}, \lean{unitCube}, \ldots); \lean{selbergInt} $\to$
    \lean{selbergIntegral}, \lean{aomotoInt} $\to$ \lean{aomotoIntegral} (decision~5).
\end{itemize}

Names flagged by the scan and kept: DLMF letters (\lean{spheroidalPs}, \lean{spheroidalQs},
\lean{discreteQHermiteI}), the suffix \lean{Int} where it means an integer index
(\lean{mathieuMeInt}, \lean{qPochhammerInt}), \lean{ODE2}, \lean{besselWronskianSeries}
and \lean{partitionCountDiff}. The suffix \lean{Big} is kept: it marks a DLMF capital letter,
the $q$-exponentials $E_q$ (\lean{qExpBig}, \lean{qSinBig}, \lean{qCosBig}) and the
alternative normalization $\Theta$ of \lean{riemannThetaBig}. \lean{Painleve.IsPIIOn}, the
Beals--Wong normalization of $\mathrm{P_{II}}$, is a bridge to \lean{IsPainleveIIOn} and is not
part of a roadmap.

\section{Decisions for the renaming round}

The following decisions were adopted on 2026-11-04.
\begin{enumerate}
  \item Parameter forms. The real functions of the parameter $m=k^2$ of the older
    \lean{Elliptic} area get the suffix \lean{OfParameter}. Roadmaps contain only the forms in
    the modulus~$k$ of Chapters~19 and~22.
  \item Complex placement per family. The whole incomplete gamma and incomplete beta families
    and \lean{multivariateGamma} go in \lean{Complex}, with names that mirror TauCeti's real
    functions (\lean{Complex.regularizedGamma}). Every other function goes at the root. The
    real-argument regularized incomplete beta function gets a bridge to TauCeti's on real
    parameters. One of the two definitions of $P$ and of $Q$ is to be kept, after a proof
    round.
  \item General areas. Chapter~1, Chapter~2 and \lean{DLMF.Core} stay out of the renaming
    round; their duplicates are merged and the names that contained ``DLMF'' are spelled
    out. They get a separate Mathlib-overlap review before any roadmap.
  \item Coulomb functions of \S33.14: \lean{coulombEnergyF}, \lean{coulombEnergyH},
    \lean{coulombEnergyS}, \lean{coulombEnergyC}, after their variables $(\epsilon,r)$;
    ``radial'' would not separate them from $F_\ell(\eta,\rho)$.
  \item Helper definitions in the foundational layer. Most are used by other modules, so
    making them private rarely works. Helpers that no foundational module needs move to the
    theorem layer; the others are spelled out as proposed. \lean{selbergIntegral} and
    \lean{aomotoIntegral} are DLMF objects (\S5.14), not helpers. The \lean{Elliptic}
    helpers \lean{piIntegrand} and \lean{piZeroPrimitive} belong to this group.
  \item Theta functions keep spelled-out numbers (\lean{jacobiThetaOneTau},
    \ldots, \lean{jacobiThetaFourQ}). Numeral subscripts are not used: Mathlib's
    \lean{jacobiTheta₂}~$z\,\tau=\sum_n e^{2\pi inz+\pi in^2\tau}$ is DLMF's
    $\theta_3(\pi z|\tau)$, the subscript counting variables, so \lean{jacobiTheta₂} for DLMF's
    $\theta_2$ would denote a different function. The bridge is recorded as a
    \lean{dlmf\_upstream} pointer with relation ``different convention''. Whether \lean{Tau}
    becomes \lean{OfTau}, as in Chapter~22, is settled with the Chapter~20 roadmap.
\end{enumerate}
\lean{hermiteFunction} is replaced by TauCeti's function: the tag 1.18.E42 is on the bridge
\lean{ofReal\_hermiteFunction}, and the two theorems of 1.18.41--1.18.43 are stated for the
cast of \lean{TauCeti.hermiteFunction}.

\section{The renaming round}

The renaming round of 2026-11-04 applied 3,221 renames: the definitions and theorems of the
table and the theorem-layer names that contain them. The names that the round left unchanged are
these.
\begin{itemize}
  \item Carlson's symmetric integrals \lean{symEllIntRF}, \lean{RC}, \lean{RD}, \lean{RG},
    \lean{RJ} and the 319 names that contain them, at first. lean-codes declares
    \lean{Carlson.carlsonRF}, \lean{Carlson.carlsonRG} and
    \lean{Carlson.TwoVariable.carlsonRC} for the same functions, and the project has bridge
    theorems between the two (\lean{symEllIntRC\_eq\_carlsonRC}). On 2026-11-04 the project
    took names of its own in the style of \lean{completeEllipticK}: \lean{symmetricEllipticRF},
    \lean{RC}, \lean{RD}, \lean{RG}, \lean{RJ} and \lean{symmetricEllipticR}, with the
    helpers \lean{homogeneousHypergeometricPolynomial} ($T_N$, 19.19.1),
    \lean{elementarySymmetric} ($E_s$, 19.19.4), \lean{quadraticTransformationP} ($p_\pm$,
    19.22.5) and \lean{duplicationLambda} ($\lambda$, 19.26.19); 407 public and 53 private
    names. Names with \lean{carlsonR} keep referring to the lean-codes functions.
  \item The two definitions of $P$ and $Q$ and the real incomplete beta bridge (decision~2).
  \item The merges of duplicates in the general areas (decision~3).
\end{itemize}
Three files that the longer names pushed past 1,500 lines were split along their DLMF
subsections: \module{DLMF.Hypergeometric.QuadraticConnections} (\S15.8(iv)),
\module{DLMF.GeneralizedHypergeometric.UnitArgumentDougall} (16.4.9--16.4.10) and
\module{DLMF.EllipticIntegrals.IncompleteInequalities} (\S19.9(ii)); the Carlson names later
split off \module{DLMF.EllipticIntegrals.SymmetricSpecialCasesRD} (\S19.20(iv)--(v)).

After the renaming round, the scan is rerun and the table regenerated. The first roadmap
is then written for a chapter whose foundational names come out of the round clean.
Candidates are Chapter~12 (parabolic cylinder functions), Chapter~22 (Jacobi elliptic
functions) and Chapter~19 (elliptic integrals).

\end{document}
